\documentclass[11pt]{amsart}
\usepackage[T1]{fontenc}
\usepackage{amsmath,amssymb,amsthm}
\usepackage[margin=1.1in]{geometry}
\usepackage{booktabs}
\usepackage{enumitem}
\usepackage{hyperref}

\newtheorem{theorem}{Theorem}[section]
\newtheorem{lemma}[theorem]{Lemma}
\newtheorem{proposition}[theorem]{Proposition}
\newtheorem{corollary}[theorem]{Corollary}
\newtheorem{conjecture}[theorem]{Conjecture}
\theoremstyle{definition}
\newtheorem{definition}[theorem]{Definition}
\newtheorem{remark}[theorem]{Remark}

\newcommand{\cD}{\mathcal D}
\newcommand{\lean}[1]{\quad{\footnotesize[Lean: \texttt{#1}]}}
\newcommand{\STL}{\mathrm{STL}}
\newcommand{\ind}[1]{[\,#1\,]}
\newcommand{\sh}{{\sharp}}

\title[The dense residual of T(6)-i]{The dense residual of the two-end test at degree six:\\
the append chain}
\author{Justin Williams}
\date{November 2026 (working draft)}

\begin{document}
\begin{abstract}
The two-end deletion test at a vertex of degree six (T(6)-i) is the first open rung of the
spine of Gallai's path-decomposition conjecture. Its open part is the \emph{dense} regime: if
$z$ is the unique even vertex of a $3$-connected graph $J$ and $U=N(z)$, a failure has
$D=K_6-(J-z)[U]\in\{\emptyset,K_2,P_3,K_3,C_4\}$. The composition method of the
terminal-matching paper reduces the test to a Lov\'asz decomposition of a base graph with one even
vertex that separates six pairs. In the dense regime that base hypothesis cannot be met for
$D\ne C_4$, because separated pairs must be non-adjacent, and for $D=C_4$ it needs a protected
Fan step.

In the dense regime the two walks of the chained walk theorem can be taken to be \emph{single
edges} of $J-z$. Every step of the construction is then an explicit append, except possibly one
protected Fan step. So the six separations can be replaced by \emph{end-path} conditions, which
ask only that the path \emph{ending} at a given vertex avoid another; this is the finer tracking
proposed in the terminal-matching paper, carried out here. These conditions do not require
non-adjacency. We prove (Theorem~\ref{thm:append}) that one Lov\'asz decomposition of the base
graph $M=J-\{zx_1,zx_2,zx_3,zy_2,zy_3,x_1y_2,x_2y_3\}$ with six end-path conditions gives the
pairing-prescribed star-terminal condition STL-P with three transitions, hence
$\STL(J,z,y_i)$ for $i=1,2,3$. One condition can be traded for a static predecessor condition
at a protected step (Theorem~\ref{thm:protected}). The first deletion test at $z$ follows from
such bases for labellings covering all six roots. No chain hypothesis is needed, and the
statement is uniform over all five dense shapes. Both theorems are formalised in Lean~4 without
\texttt{sorry} and with the standard axioms only.

The append form is reversible (Proposition~\ref{prop:converse}). A base with the append
conditions exists if and only if $J$ has a star-terminal Lov\'asz decomposition with
transitions $x_i\,z\,y_i$ in which, in addition, $x_1y_2$ is terminal at $y_2$ and $x_2y_3$ is
terminal at $y_3$. So the base conjecture~\ref{conj:base} is a \emph{strengthening} of STL-P
by two prescribed terminal edges, relaxed at $y_2$ by the protected step. It is a uniform
re-encoding of the dense residual, not a reduction to a weaker statement.

In the computed instances a base exists for every root except $46$ roots in eight graphs, where
the search hit its budget. This holds for all five shapes, including random $3$-connected graphs
carrying a subset of the residual conditions of a smallest failure: minimum degree six,
odd-neighbour counts, contact sizes, closed twins for $C_4$ and a full exterior vertex for
$\emptyset$. Within the same budget, the six-separation base of the chained walk theorem was
never found in these graphs. \emph{The existence of the base is not proved}, and the dense
residual of T(6)-i, $C_4$ included, is not closed.
\end{abstract}
\maketitle

\section{Introduction}\label{sec:intro}

All graphs are finite and simple. A \emph{Lov\'asz decomposition} of a graph with at most one
vertex of positive even degree is a decomposition into paths in which every vertex of odd degree
is an end of exactly one path and every vertex of even degree is an end of none.

\subsection*{The test}
Let $z$ be the only vertex of even degree of a $3$-connected graph $J$, with $d(z)=6$, and let
$B=J-z$ and $U=N(z)$. The first deletion test $\mathrm{Test1}(J,z)$ asks that, for every
$r\in U$, $B$ have a decomposition into at most $(|J|-1)/2$ paths with at least two ends at
$r$. By the spine~\cite{SP} (Lean \texttt{gallai\_iff\_spineTests'}), Gallai's conjecture is
equivalent to this test together with a second one. At $d(z)\le4$ the test is known, and at
$d(z)=6$ it is the open case T(6)-i~\cite{RS,INC,MR}.

A decomposition of $J$ is \emph{star-terminal} at $z$ if every path through $z$ has the form
$x\,z\,y\cdots$ with $x$ an end of the path. The neighbour $y$ is then on a \emph{long side},
unless the path is $x\,z\,y$. $\STL(J,z,r)$ asks for a star-terminal Lov\'asz decomposition of
$J$ in which $r$ is on a long side, or which has a path of length two through $z$. By
\texttt{test1\_of\_stl}~\cite{STL}, $\STL(J,z,r)$ for every $r\in U$ (with every neighbour of
degree at least three) gives $\mathrm{Test1}(J,z)$.

\subsection*{The residual}
Write $C=B[U]$ (the \emph{core}) and $D=K_6-C$. By the root-shape theorem~\cite{RS}, a failure
of the test at a root $r$ \emph{with a neighbour in $C$} has $D\in\{\emptyset,K_2,P_3,K_3,C_4\}$,
with $r$ restricted further ($r$ the middle vertex of $P_3$, a vertex of $K_3$, a vertex of the
$C_4$ whose two cycle-neighbours are adjacent twins throughout $B$). A root isolated in $C$ has
$D\supseteq K_{1,5}$, outside the five shapes. Such roots are excluded by \cite[Thm~144]{MR}:
failing roots may be taken with $d(r)\ge8$ and at least three neighbours among the other even
vertices. (Audit A6, F16--F18, removes the unrefereed input of its proof.) Many further conditions hold in a smallest failure (STATUS~\S3.5(a);
reader notes on the six-even program): $\delta(B)\ge6$, $\delta(B[O])\ge4$ for the set $O$ of
odd vertices, every even site has at least three odd neighbours, $|N(x)\cap U|\in\{0,3,4,5,6\}$
for $x\in O$, and $D=\emptyset$ only with an exterior vertex adjacent to all of $U$. Among the
shapes, $D\in\{K_2,P_3,K_3\}$ has no shape-specific result, $C_4$ is open for twins of degree at
least eight, and $\emptyset$ for full exterior vertices.

\subsection*{The composition method and its obstruction}
For a pairing $F=\{x_1y_1,x_2y_2,x_3y_3\}$ of $U$ the terminal-matching paper~\cite{TM}
reduces $\STL$ at $y_1,y_2,y_3$ to a Lov\'asz decomposition of $(J-z)+F$ in which every
$x_iy_i$ is terminal at $x_i$, on three different paths. Its chained walk theorem
(\cite[Thm~8.2]{TM}, Lean \texttt{ptm3\_of\_base}) does three things.
\begin{itemize}
\item It replaces $x_1y_1$ by a leaf $\ell_1$ at $y_1$.
\item It walks the even vertex twice, along a path $P_2$ from $y_3$ to $x_2$ and a path $P_1$
from $y_2$ to $x_1$, restoring the paths edge by edge with Fan's lemma.
\item It asks for a Lov\'asz decomposition of the base graph separating six pairs:
$\{x_1,\ell_1\}$, $\{x_2,y_2\}$, $\{x_3,y_3\}$, $\{\ell_1,y_2\}$, $\{\ell_1,y_3\}$,
$\{y_2,y_3\}$.
\end{itemize}
In the dense regime this fails twice over (synthesis~\S22; connections~\S Z.4):
\begin{itemize}
\item every edge lies on a path, so separated pairs must be non-adjacent in the base graph, and
this forces the path $x_2y_2y_3x_3$ into $D$, which only $C_4$ contains;
\item for $C_4$ the vertex $y_2$ is adjacent to $x_3$, so the first step of the second walk is a
Fan step at equality, with the endless $x_3$ as a blank neighbour of the target. It needs the
protected form \texttt{walk\_step\_protected} of Fan's lemma~\cite{TF}.
\end{itemize}

\subsection*{What this paper does}
\emph{Chains are free in the dense regime.} If $x_1y_2$ and $x_2y_3$ are edges of $J$, we may take
$P_2=y_3x_2$ and $P_1=y_2x_1$. Such a labelling exists for every root, for every graph $D$ of
maximum degree at most two, so for all five shapes (Lemma~\ref{lem:labelling}). With walks of
length one, every step of the construction is explicit:
\begin{itemize}
\item a leaf at $y_3$;
\item an \emph{append} $x_2\to y_3$;
\item a leaf at $y_2$;
\item the step $x_1\to y_2$, which is an append or one protected Fan step;
\item three appends at the leaves.
\end{itemize}
What must avoid $x_i$ is only the path \emph{ending} at $y_i$. This finer tracking is the one
proposed in \cite[\S9(1)]{TM}. With walks of length one, that path can be named at each step. So
the base hypothesis becomes a list of end-path conditions (Definition~\ref{def:base}).
\begin{itemize}
\item \textbf{Theorem~\ref{thm:append}} (append chain). A base with the end-path conditions
(E1)--(E5) and (E7) gives a star-terminal Lov\'asz decomposition of $J$ with transitions
$x_i\,z\,y_i$, so $\STL(J,z,y_i)$ for $i=1,2,3$.
\lean{stl\_of\_appendBase}
\item \textbf{Proposition~\ref{prop:converse}} (reversibility). Conversely, such a base exists if
and only if $J$ has a star-terminal Lov\'asz decomposition with transitions $x_i\,z\,y_i$ in
which $x_1y_2$ is terminal at $y_2$ and $x_2y_3$ is terminal at $y_3$.
\item \textbf{Theorem~\ref{thm:protected}} (one protected step). The same with (E7) replaced by
(E6). This is the static predecessor condition of \texttt{walk\_step\_protected} at the step
$x_1\to y_2$, transported back to the base decomposition through the three appends
(Lemma~\ref{lem:transport}).
\lean{star\_of\_protectedBase}
\item \textbf{Corollary~\ref{cor:test1}}. $\mathrm{Test1}(J,z)$ holds as soon as every
$r\in U$ admits a labelling with $r\in\{y_1,y_2,y_3\}$ and a base.
\lean{test1\_of\_appendBases}
\item \textbf{Data} (E28). For all five shapes, in the computed instances, every root has a
labelling with a base, except $46$ roots in eight graphs where every search hit the node budget
(Section~\ref{sec:comp}). The instances include $3$-connected graphs carrying a subset of the
residual conditions above. In those graphs the six-separation base of \cite[Thm~8.2]{TM} was
never found within the search budget.
\end{itemize}
The remaining hypothesis is Conjecture~\ref{conj:base}, the existence of a base. By
Proposition~\ref{prop:converse} it is, up to the protected escape (E6), the pairing-prescribed
star-terminal condition STL-P of \cite{STL} (synthesis~\S20) for some labelled pairing, together
with two further prescribed terminal edges: on the alternating path $y_1x_1y_2x_2y_3x_3$ of
$(J-z)+F$, every edge is terminal at its forward end, and the three edges of $F$ lie on
different paths (Remark~\ref{rem:ptm}). So it is a \emph{strengthening} of
STL-P, which in turn is a strengthening of the condition $\STL$ that gives the test. It is
\emph{not proved}. In particular the $C_4$ shape of T(6)-i is not closed by this paper.

\subsection*{Novelty}
Section~\ref{sec:novelty} compares the result with the known T(6)-i reductions. In short, none
of the known closures (six-even-incidence, the four-exterior and five-odd papers, the
six-clique theorem, the twin results, the degree-ten note) gives the dense shapes in general.
Several known reductions already have one even vertex: $\STL$ via \texttt{test1\_of\_stl},
STL-P and its terminal-matching form PTM$_3$, and \cite[Thm~8.2]{TM}. The previous composition
result \cite[Cor.~8.4]{TM} covers only the far-apart case, which was already known
\cite[Rem.~8.5]{TM}. The end-path tracking that removes the obstruction of \cite[\S9(1)]{TM}
was proposed there. What is new here:
\begin{itemize}
\item the chain design with walks of length one on the cross edges $x_1y_2$, $x_2y_3$, which
exists for every root in all five dense shapes (Lemma~\ref{lem:labelling});
\item the exact end-path conditions for it, and the transport of the static predecessor
condition through the appends (Theorem~\ref{thm:protected});
\item the converse (Proposition~\ref{prop:converse}), which identifies the base conjecture as
STL-P with two further prescribed terminal edges;
\item the Lean formalisation of the construction.
\end{itemize}
No new case of T(6)-i is closed, and the base conjecture is not shown to be easier than STL-P.

\section{The base and the main theorems}\label{sec:main}

Throughout, $J$ is a graph, $z$ a vertex with $N(z)=\{x_1,y_1,x_2,y_2,x_3,y_3\}$ (seven distinct
vertices), and $z$ is the only vertex of even degree of $J$. In a decomposition in which every
vertex has at most one end, $P_v$ denotes the path ending at $v$. For a path $S$ through an edge
$py$, the \emph{end of $S$ on $p$'s side} is the end $e$ of $S$ such that, reading $S$ from $e$,
the vertex $p$ comes immediately before $y$. In Fan's terminology~\cite{Fan05,TF}, $p$ is then
the predecessor of $y$ read from $e$, and we write $\mathrm{pred}_y(e)=p$.

\begin{definition}[base graph, append base]\label{def:base}
Suppose $x_1y_2,x_2y_3\in E(J)$. The \emph{base graph} is
\[
M=J-\{zx_1,zx_2,zx_3,zy_2,zy_3,x_1y_2,x_2y_3\}.
\]
In $M$ the vertex $z$ is a leaf at $y_1$, $x_3$ has even degree, and every other vertex has odd
degree. (The degrees of $x_1,x_2,y_2,y_3$ drop by two, that of $x_3$ by one.) A Lov\'asz
decomposition $\cD_0$ of $M$ has one end at every vertex except $x_3$, and none at $x_3$. It is an
\emph{append base} if
\begin{enumerate}[label=\textup{(E\arabic*)}]
\item the path through $z$ (that is, $P_z$) avoids $x_1$, and its other end is neither $y_2$
nor $y_3$;
\item $P_{y_2}$ avoids $x_2$;
\item $P_{y_3}$ avoids $x_3$;
\item $P_{y_2}\ne P_{y_3}$, that is, no path has ends $y_2$ and $y_3$;
\item $P_{x_2}$ avoids $y_3$;
\end{enumerate}
and one of the following holds:
\begin{enumerate}[label=\textup{(E\arabic*)}]\setcounter{enumi}{5}
\item for $p\in\{x_2,x_3\}$, no $w\in N_J(y_2)\setminus\{z,x_2\}$ has $\mathrm{pred}_{y_2}(w)=p$
in $\cD_0$, and if $y_2y_3\in E(J)$, then also $\mathrm{pred}_{y_2}(x_2)\ne p$ (the static
predecessor condition);
\item $P_{x_1}$ avoids $y_2$.
\end{enumerate}
\end{definition}

In Lean, \texttt{LovaszBase} is Definition~\ref{def:base} verbatim: a decomposition
$\cD_0$ of $M$ (\texttt{baseM}) with ends equal to degrees modulo two, and (E1)--(E7) on $V$. The
degree parities of $M$ are \texttt{baseM\_parity}. \texttt{AppendBase} is the same in the split
form of Section~\ref{sec:proof}: $M$ together with two isolated vertices $\ell_2,\ell_3$, with
the explicit end profile. \texttt{appendBase\_of\_lovaszBase} pushes the first into the second.

\begin{theorem}[the append chain]\label{thm:append}
If $M$ has a Lov\'asz decomposition satisfying \textup{(E1)}--\textup{(E5)} and
\textup{(E7)}, then $J$ has a star-terminal Lov\'asz decomposition in which the three paths
through $z$ are $x_i\,z\,y_i\cdots$ with $x_i$ an end ($i=1,2,3$). In particular
$\STL(J,z,y_i)$ holds for $i=1,2,3$.\lean{stl\_of\_lovaszBase}
\end{theorem}

\begin{theorem}[one protected step]\label{thm:protected}
The conclusion of Theorem~\ref{thm:append} also holds if \textup{(E7)} is replaced by
\textup{(E6)}.
\end{theorem}

The Lean statements (\texttt{star\_of\_appendBase}, \texttt{star\_of\_protectedBase} and their
corollaries) cover both theorems with (E6) or (E7). They record what $\STL$ needs: a Lov\'asz
decomposition of $J$ that is star-terminal at $z$, in which each $y_i$ is on a long side or $z$
lies on a path of length two. The transitions $x_i\,z\,y_i$ are built in the proof
(\texttt{split\_finish}) but are not part of the Lean conclusion.

\begin{corollary}\label{cor:test1}
Let $z$ be the only even vertex of $J$, $d(z)=6$, and let every neighbour of $z$ have degree at
least three.
\begin{enumerate}[label=\textup{(\alph*)}]
\item If $r\in N(z)$ and some labelling $(x_1,y_1,x_2,y_2,x_3,y_3)$ of $N(z)$ with
$r\in\{y_1,y_2,y_3\}$ has an append base, then $B=J-z$ has a decomposition into at most
$(|J|-1)/2$ paths with at least two ends at $r$.
\item If this holds for every $r\in N(z)$, then $\mathrm{Test1}(J,z)$ holds.
\lean{test1\_of\_lovaszBases}
\end{enumerate}
\end{corollary}

\begin{proof}
By Theorems~\ref{thm:append} and~\ref{thm:protected}, $\STL(J,z,r)$ holds. Part (b) is then
\texttt{test1\_of\_stl}~\cite{STL}. Its proof uses $\STL(J,z,r)$ only at the root $r$ under
test, which gives (a). Part (a) is not stated separately in Lean.
\end{proof}

\begin{lemma}[labellings exist]\label{lem:labelling}
Let $D=K_6-B[U]$ have maximum degree at most two, which holds for all five residual shapes. Then
for every $r\in U$ there is a labelling with $y_2=r$ and $x_1y_2,x_2y_3\in E(B)$. The same holds
with $y_3=r$, and with $y_1=r$.
\end{lemma}

\begin{proof}
For $y_2=r$: $r$ has at most two $D$-neighbours among the other five vertices, so we can choose
$x_1$ with $rx_1\notin D$. A graph on the remaining four vertices with maximum degree at most two
is not $K_4$, so it has a non-edge $x_2y_3$. For $y_3=r$, choose $x_2$ with $rx_2\notin D$ and
then a non-edge $x_1y_2$ among the remaining four vertices. For $y_1=r$, we need two disjoint edges $x_1y_2,x_2y_3$ of
$B$ among the other five vertices. On five vertices $D$ has at most five edges, so $B$ has at least
five there. A graph with no two disjoint edges is a star or a triangle, with at most four edges on
five vertices, so two disjoint edges exist.
\end{proof}

So in the dense residual the chain hypothesis of \cite[Def.~8.1]{TM} and Conjecture~8.6 there
are not used. (The length-one chain violates \cite[Def.~8.1(ii)]{TM} whenever $y_2x_3\in E(B)$,
and Conjecture~8.6 there concerns far-apart neighbours only.) For the $C_4$ shape and a root $r$ on the cycle, the labelling that runs
$x_2\,y_2\,y_3\,x_3$ along the cycle, with $\{x_1,y_1\}$ the two universal vertices, has
$x_1y_2$ (a universal vertex) and $x_2y_3$ (a diagonal) in $B$. This is the configuration of
\cite[\S9(1)]{TM}. There $y_2x_3$ is the other diagonal, so the step $x_1\to y_2$ has the blank
$x_3$, and it is either protected, by (E6), or an append, by (E7).

\begin{proposition}[the append form is reversible]\label{prop:converse}
Let $z$ be the only even vertex of $J$, $N(z)=\{x_1,y_1,x_2,y_2,x_3,y_3\}$, and
$x_1y_2,x_2y_3\in E(J)$. The following are equivalent.
\begin{enumerate}[label=\textup{(\alph*)}]
\item $M$ has a Lov\'asz decomposition satisfying \textup{(E1)}--\textup{(E5)} and
\textup{(E7)}.
\item $J$ has a Lov\'asz decomposition $\cD$ such that every path through $z$ has the form
$S_i=x_i\,z\,y_i\cdots$ with $x_i$ an end ($i=1,2,3$), the path ending at $y_2$ ends with the edge
$x_1y_2$, and the path ending at $y_3$ ends with the edge $x_2y_3$.
\end{enumerate}
\end{proposition}

The proof is in Section~\ref{sec:proof}; it is not formalised. The auditor of A26 found the
statement and confirmed it by exact search on all $14{,}144$ labellings with chain edges of the
$28$ dense graphs of order $9$. Both sides hold for $14{,}128$ labellings and fail for $16$, with no
disagreement.

\begin{remark}[the terminal-matching form]\label{rem:ptm}
Let $F=\{x_1y_1,x_2y_2,x_3y_3\}$ and $G=(J-z)+F$ (a multigraph if some $x_iy_i\in E(B)$). Replacing
$x_i\,z\,y_i$ by the edge $x_iy_i$ turns (b) into the following statement. $G$ has a Lov\'asz
decomposition in which the three edges of $F$ lie on different paths and every edge of the
alternating path $y_1\,x_1\,y_2\,x_2\,y_3\,x_3$ is terminal at its forward end: $y_1x_1$ at $x_1$,
$x_1y_2$ at $y_2$, $y_2x_2$ at $x_2$, $x_2y_3$ at $y_3$ and $y_3x_3$ at $x_3$. Without the two cross
edges this is PTM$_3$ for $F$ with prescribed ends, that is, STL-P for the oriented pairing
\cite{STL,TM} (synthesis~\S20). So Theorem~\ref{thm:append} derives STL-P from a
\emph{stronger} statement, and Conjecture~\ref{conj:base} restricted to (E7) is that stronger
statement. The (E6) branch relaxes the requirement at $y_2$; it is rarely needed in the data
(Section~\ref{sec:comp}).
\end{remark}

\section{Proofs}\label{sec:proof}

\subsection*{The split graph}
Split $z$ into three vertices of degree two:
\begin{itemize}
\item $z$ itself, adjacent to $x_1,y_1$;
\item $\ell_2$, adjacent to $x_2,y_2$;
\item $\ell_3$, adjacent to $x_3,y_3$.
\end{itemize}
Call the result $J^\sh$ (Lean \texttt{splitG}, on $V\oplus\mathrm{Fin}\,2$). Identifying
$z,\ell_2,\ell_3$ maps $E(J^\sh)$ bijectively onto $E(J)$ (\texttt{collapse\_hE},
\texttt{collapse\_hEinj}, \texttt{collapse\_hcov}). Let
\[
L=(\ell_3y_3,\ x_2y_3,\ \ell_2y_2,\ x_1y_2,\ zx_1,\ \ell_2x_2,\ \ell_3x_3).
\]
Then $J^\sh-L$ is $M$ plus the isolated vertices $\ell_2,\ell_3$. The construction restores the
seven edges of $L$ in order; the graphs in between are $J^\sh$ minus a final segment of $L$
(\texttt{delList}, \texttt{spl}).

\subsection*{Appends}
Let $\cD$ decompose a graph $K$, let $tu\notin E(K)$, and let $P\in\cD$ have an end at $t$ and
avoid $u$. Appending $tu$ at $t$ gives a decomposition $\cD'$ of $K+tu$. The new path $P'$ starts
at $u$ and, read from $u$, is $u$ followed by $P$ read from $t$. Every other path of $\cD'$ has
the ends and the vertex sequence of a path of $\cD$, and the ends move from $t$ to $u$
(\texttt{append\_full}). We say $P'$ \emph{extends} $P$ at $t$ by $u$. An extended path has the
edges of $P$ and the edge $tu$ (\texttt{ExtP.edges\_mem}).

\begin{lemma}[predecessors through an append]\label{lem:transport}
In this situation let $w\ne u$ and $p\ne u$. If $\mathrm{pred}_w(x)=p$ in $\cD'$, then
$\mathrm{pred}_w(x)=p$ in $\cD$, or $x=u$ and $\mathrm{pred}_w(t)=p$ in $\cD$. Moreover
$\mathrm{pred}_w(w)$ is never defined.\lean{fanPred\_append, fanPred\_self}
\end{lemma}

\begin{proof}
A path of $\cD'$ other than $P'$ has the ends and the vertex sequence of a path of $\cD$. On
$P'$, read from $u$, the pair $(p,w)$ occurs inside $P$ read from $t$, since $p,w\ne u$. Read
from its other end, which is an end of $P$, the pair occurs inside $P$ read from that end. A path
read from $w$ starts at $w$, so nothing precedes $w$.
\end{proof}

\subsection*{Proof of Theorem~\ref{thm:append}}
Let $\cD_0$ satisfy (E1)--(E5) and (E7), viewed as a decomposition of $J^\sh-L$. The leaves
$\ell_2,\ell_3$ are isolated there, so no path contains them.

\emph{Step 1 (leaf at $y_3$).} Append $\ell_3$ to $P_{y_3}$. This is possible because $\ell_3$
is isolated. Now $y_3$ is endless and $\ell_3$ has an end.

\emph{Step 2 (the first walk, $x_2\to y_3$).} The path ending at $x_2$ is a path of $\cD_0$ with
an end at $x_2$. It cannot be the extension of $P_{y_3}$, since $P_{y_3}$ contains $y_3$ and
(E5) would be violated. So it avoids $y_3$ by (E5), and we append $y_3$ to it. Now $x_2$ is
endless and $y_3$ has an end again.

\emph{Step 3 (leaf at $y_2$).} $\ell_2$ is isolated, so we can append $\ell_2$ to the path ending
at $y_2$. That path is $P_{y_2}$ of $\cD_0$. It is not the extension of $P_{x_2}$, since then
$P_{x_2}$ would end at $y_2$ and contain $x_2$, against (E2). It is not the extension of
$P_{y_3}$, by (E4).

Call the result $\cD_3$. Every path of $\cD_3$ is a path of $\cD_0$ (same ends and vertex
sequence) or a path $Q$ of $\cD_0$ extended once:
\begin{itemize}
\item at $y_3$ by $\ell_3$, with $Q=P_{y_3}$;
\item at $x_2$ by $y_3$, with $Q=P_{x_2}$;
\item at $y_2$ by $\ell_2$, with $Q=P_{y_2}$.
\end{itemize}
This is \texttt{split\_abc}. Its end counts are those of $\cD_0$, except that
$x_2,y_2$ have no end and $\ell_2,\ell_3$ have one.

\emph{Leaf conditions in $\cD_3$} (\texttt{leafOK\_three}):
\begin{itemize}
\item the path through $z$ avoids $x_1$, $\ell_2$, $\ell_3$;
\item the path through $\ell_2$ avoids $x_2$ and $\ell_3$;
\item the path through $\ell_3$ avoids $x_3$.
\end{itemize}
They follow from the classification together with (E1)--(E3). For example, an extension of
$P_{y_3}$ through $z$ would make $P_z$ end at $y_3$, against (E1). The path through $\ell_3$ is
$P_{y_3}$ plus $\ell_3$, which avoids $x_3$ by (E3).

\emph{Step 4 (the second walk, $x_1\to y_2$).} The path $R$ of $\cD_3$ ending at $x_1$ avoids
$y_2$:
\begin{itemize}
\item if $R$ comes from a path $Q$ of $\cD_0$, unchanged or extended at $y_3$ or at $x_2$, then
$Q$ ends at $x_1$, and (E7) gives $y_2\notin Q$; the added vertices are not $y_2$;
\item it cannot be an extension of $P_{y_2}$, since then $P_{y_2}$ would end at $x_1$ and
contain $y_2$, against (E7).
\end{itemize}
Append $y_2$ to $R$. Every new path lies in an old path together with $y_2$, so the leaf
conditions survive (\texttt{LeafOK.of\_cover}); none of $z,x_1,x_2,x_3$ is $y_2$. Now
$x_1,x_2,x_3$ have no ends, the three leaves $z,\ell_2,\ell_3$ have one each, and every other
vertex has one.

\emph{Steps 5--7.}
\begin{itemize}
\item Append $x_1$ to the path ending at $z$. It avoids $x_1$ by the leaf conditions.
\item Append $x_2$ to the path ending at $\ell_2$. It is not the path just extended, since that
path contains $z$ and $\ell_2$ is not on it; so it is an old path, which avoids $x_2$.
\item Append $x_3$ to the path ending at $\ell_3$, likewise.
\end{itemize}
In the resulting decomposition $\cD_7$ of $J^\sh$, each of $z,\ell_2,\ell_3$ has degree two and
no end. Each of them lies on exactly one path $S_1,S_2,S_3$, which starts $x_i\,\ell_i\,y_i$ with
$\ell_1=z$, and each $S_i$ contains only one of the three. This uses that a leaf of degree one
lies on the path that ends at it.

\emph{Step 8 (identify).} Mapping $\ell_2,\ell_3$ to $z$ is injective on each path, and maps
edges bijectively, so it pushes $\cD_7$ to a decomposition $\cD$ of $J$
(\texttt{PathDecomp.pushOn}). The vertex $z$ has no end and every other vertex has one, so $\cD$
is a Lov\'asz decomposition of $J$. A path through $z$ is the image of some $S_i$, which starts
$x_i\,z\,y_i$, so $\cD$ is star-terminal. If $S_i$ ends at $y_i$, then before step $5$ the leaf
path was the single edge $\ell_iy_i$, because $y_i$ is the only neighbour of $\ell_i$ there. So
the image is $x_i\,z\,y_i$, a path of length two through $z$. Otherwise $y_i$ is on a long side.
This is \texttt{split\_finish}; \texttt{stl\_of\_appendBase} adds the degree parities.
\qed

\subsection*{Proof of Theorem~\ref{thm:protected}}
Steps 1--3 and 5--8 are as before. Step 4 is now Fan's transfer at the target $y_2$, restoring
$y_2x_1$, in the protected form \texttt{walk\_step\_protected} with protected set $\{x_2,x_3\}$.
Its hypotheses are the following.
\begin{itemize}
\item Every neighbour of $y_2$ other than $x_1$ and outside $\{x_2,x_3\}$ has an end in
$\cD_3$. Indeed every vertex except $x_2,x_3,y_2$ has one, and $\ell_2$ has one.
\item $x_1$ has an end.
\item \emph{The static predecessor condition:} no vertex of $\{x_1\}\cup N(y_2)$ has
$\mathrm{pred}_{y_2}\in\{x_2,x_3\}$ in $\cD_3$.
\end{itemize}
Apply Lemma~\ref{lem:transport} three times, using $y_2\notin\{\ell_3,y_3,\ell_2\}$ and
$p\notin\{\ell_2,y_3,\ell_3\}$. Then $\mathrm{pred}_{y_2}(x)=p$ in $\cD_3$ implies one of:
\begin{itemize}
\item $\mathrm{pred}_{y_2}(x)=p$ in $\cD_0$;
\item $x=\ell_3$ and $\mathrm{pred}_{y_2}(y_3)=p$ in $\cD_0$;
\item $x=y_3$ and $\mathrm{pred}_{y_2}(x_2)=p$ in $\cD_0$.
\end{itemize}
The step at $y_2$ itself contributes nothing, by $\mathrm{pred}_{y_2}(y_2)$ undefined. Now
$\ell_3$ is not adjacent to $y_2$. The vertex $x=x_2$ has no end in $\cD_3$, so it cannot be read
from. The vertex $x=\ell_2$ is isolated in $\cD_0$. Every other neighbour $x$ of $y_2$ in
$J^\sh-\{x_1y_2,zx_1,\ell_2x_2,\ell_3x_3\}$ is a vertex $w\in N_J(y_2)\setminus\{z,x_2\}$,
because $zy_2$ is not an edge of $J^\sh$. So (E6) is exactly what excludes the remaining cases.
The step gives a decomposition with the end counts of step 4 in which every path lies in an old
path together with $y_2$, and the proof finishes as before.\qed

\begin{remark}[a redundant clause]\label{rem:e6}
The clause $w=y_3$ of (E6) is not needed (A26, \S1.2). For $x=y_3$, the outcome
``$\mathrm{pred}_{y_2}(y_3)=p$ in $\cD_0$'' can only come, through step~2 (an append at $x_2$
with $u=y_3$), from $\mathrm{pred}_{y_2}(y_3)=p$ in $\cD_1$. That would read a path of $\cD_1$
from the end $y_3$, but $y_3$ has no end in $\cD_1$. So Theorem~\ref{thm:protected} holds with
(E6) restricted to $w\ne y_3$. The Lean statement keeps the clause.
\end{remark}

\subsection*{Proof of Proposition~\ref{prop:converse}}
\emph{(a)$\Rightarrow$(b).} Run the construction of Theorem~\ref{thm:append}. Step~2 gives $y_3$
an end with terminal edge $x_2y_3$, and step~4 gives $y_2$ an end with terminal edge $x_1y_2$.
Every later append is made at another end ($x_1$ in step~4, then $z,\ell_2,\ell_3$), so both
terminal edges survive into $\cD_7$ and, by step~8, into $\cD$.

\emph{(b)$\Rightarrow$(a).} A path through $z$ uses two edges at $z$, so the three paths $S_i$
are distinct and contain every edge at $z$. Putting $\ell_i$ in place of $z$ on $S_i$ (with
$\ell_1=z$) lifts $\cD$ to a decomposition of $J^\sh$ with the same ends. Let $T$ and $T'$ be
the paths ending at $y_2$ and $y_3$. Every vertex other than $z$ has exactly one end.
\begin{enumerate}[label=(\roman*)]
\item $T\notin\{S_1,S_2\}$ and $T'\notin\{S_2,S_3\}$. Indeed $S_1$ starts at $x_1$, so it does
not end with $x_1y_2$, and $S_2$ starts at $x_2$, so it does not end with $x_2y_3$. If $T=S_2$, then
$y_2$, the third vertex of $S_2$, is an end, so $S_2=x_2zy_2$ ends with $zy_2$. The case $T'=S_3$
is the same.
\item $S_2$ and $S_3$ have length at least three. If $S_2=x_2zy_2$, then $y_2$ is an end of
$S_2\ne T$. The same holds for $S_3$.
\item $T$ and $T'$ have length at least two, because $x_1$ and $x_2$ already end on $S_1$ and $S_2$.
\end{enumerate}
Delete, in this order, $x_1\ell_1,x_2\ell_2,x_3\ell_3$, then $x_1y_2$, $x_2y_3$, $\ell_2y_2$ and
$\ell_3y_3$. Each is a terminal edge of its path at that moment: the first three start the lifted
$S_i$, the next two end $T$ and $T'$, and the last two start what remains of $S_2$ and $S_3$. Each
of these paths has length at least two at that moment, so every deletion leaves a decomposition:
\begin{itemize}
\item for $S_i$, by (ii) and $|S_1|\ge2$;
\item for $T$, by (iii), and if $T=S_3$ because $\ell_3,y_3,x_1,y_2$ are distinct;
\item for $T'$, by (iii); if $T'=S_1$ because $z,y_1,x_2,y_3$ are distinct, and if $T'=T$
because $y_3,x_2,x_1$ are distinct;
\item for $S_2$ by (ii), since $S_2\notin\{T,T'\}$; for $S_3$ by (ii), and if $T=S_3$ because
what remains, $\ell_3y_3\cdots x_1$, has length at least two.
\end{itemize}
The result $\cD_0$ decomposes $J^\sh-L$, that is, $M$ with $\ell_2,\ell_3$ isolated. Each deletion
moves one end. In total $x_3$ loses its end, $z$ gains one, and every other vertex keeps one, so
$\cD_0$ is a Lov\'asz decomposition of $M$. The conditions hold as follows.
\begin{itemize}
\item[(E1)] $P_z$ is $S_1-x_1$, cut back from $y_3$ to $x_2$ if $T'=S_1$. It avoids $x_1$. Its far
end is $x_2$ if $T'=S_1$, and otherwise the far end of $S_1$, which is neither $y_2$ (as
$T\ne S_1$) nor $y_3$ (as $T'\ne S_1$).
\item[(E2)--(E4)] $P_{y_2}$ is $S_2$ without its first two vertices, and $P_{y_3}$ is $S_3$
without its first two vertices, shortened at the far end if $T=S_3$. They are different paths,
and they avoid $x_2$ and $x_3$ respectively, because $x_i$ occurs on $S_i$ only as its first
vertex.
\item[(E5)] $P_{x_2}$ is what remains of $T'$, which lost $y_3$, its only occurrence.
\item[(E7)] $P_{x_1}$ is what remains of $T$, which lost $y_2$, its only occurrence.\qed
\end{itemize}

Under (E6) the step $x_1\to y_2$ is a Fan transfer, and the terminal edge at $y_2$ need not be
$x_1y_2$. So Proposition~\ref{prop:converse} concerns the (E7) branch only.

\begin{remark}[why the obstruction disappears]
In the chained walk theorem a separation $\{x_2,y_2\}$ was needed for \emph{every} path through
$y_2$, because after a Fan walk one only knows that each new path lies in some old path plus the
targets. When the walks are appends, the path that will be extended at $y_2$ is known: it is
$P_{y_2}$. So only $P_{y_2}$ must avoid $x_2$. The edge $x_2y_2$, if present, may lie on any
other path. This is the end-path tracking asked for in \cite[\S9(1)]{TM}, in the form where it is
exact.
\end{remark}

\section{End-path tracking for general walks}\label{sec:endpath}

For walks of length greater than one, Fan's transfers reroute paths, and the path ending at a
vertex can change. It changes in a controlled way, which we record for later use; it is not
needed for the theorems above.

\begin{lemma}[Fan with end paths]\label{lem:fanend}
Let $\cD$ decompose $K-tu$, and let $\cD'$ arise from Fan's transfer along a chain
$t=x_0,x_1,\dots,x_k$ (\texttt{fan\_exec}). Let $v\ne u$ and $v\notin\{x_1,\dots,x_k\}$; this
holds in particular if $v\notin N_{K-tu}(u)$. Then every path of $\cD'$ with an end at $v$ has
its vertex set inside $V(Q)\cup\{u\}$ for a path $Q$ of $\cD$ with an end at $v$.
\end{lemma}

\begin{proof}
A reroute along $x_i$ replaces the path $x_i\cdots x_{i+1}\,u\cdots e$ by
$x_{i+1}\cdots x_i\,u\cdots e$. It keeps the vertex set and the end $e$, and moves the end $x_i$
to $x_{i+1}$; all other paths are unchanged. The final append adds $u$ to a path at its end
$x_k$ and keeps its other end. So an end at $v\notin\{x_1,\dots,x_k,u\}$ stays on a path with the
same vertex set, plus possibly $u$. The chain vertices $x_1,\dots,x_k$ are predecessors of $u$,
hence neighbours of $u$.
\end{proof}

Applied along a walk $w_0\cdots w_k$, the end paths at every vertex outside
$T\cup N(T)$, where $T=\{w_0,\dots,w_{k-1}\}$ is the set of targets, are tracked up to $T$.
Inserting this into the proof of \cite[Thm~8.2]{TM} gives the following variant. It is not
formalised. We write $T_2=\{q_0,\dots,q_{s-1}\}$ for the targets of $P_2=q_0\cdots q_s$ and use
the notation of that proof ($Q_1$ the path through $\ell_1$, $Q_3$ the path ending at $y_3$, $Q_2$
the path ending at $y_2$ after the first walk).

\begin{proposition}[Thm~8.2 with end paths]\label{prop:ep}
Let $P_1,P_2$ be an admissible chain \cite[Def.~8.1]{TM} with $y_2\notin N(\{q_0,\dots,q_{s-1}\})$.
Suppose the base graph has a Lov\'asz decomposition in which:
\begin{itemize}
\item the path through $\ell_1$ avoids $x_1$ and does not end at $y_2$ or $y_3$;
\item $P_{y_2}$ avoids $x_2$;
\item $P_{y_3}$ avoids $x_3$;
\item $P_{y_2}\ne P_{y_3}$.
\end{itemize}
Then the conclusion of \cite[Thm~8.2]{TM} holds.
\end{proposition}

\begin{proof}
The proof of \cite[Thm~8.2]{TM} uses the six separations at the following places. Each use is
replaced as shown; the rest of that proof is unchanged.
\begin{center}\footnotesize
\begin{tabular}{@{}lp{0.36\textwidth}p{0.38\textwidth}@{}}
\toprule
separation & use in \cite[Thm~8.2]{TM} & replacement\\
\midrule
$\{x_1,\ell_1\}$ & Step 4: the $\ell_1$-path avoids $x_1$ & $Q_1$ avoids $x_1$ (kept)\\
$\{x_3,y_3\}$ & Step 1: $Q_3$ avoids $x_3$ & $P_{y_3}$ avoids $x_3$\\
$\{\ell_1,y_3\}$ & Step 1: $Q_3\ne Q_1$ & $Q_1$ does not end at $y_3$\\
$\{y_2,y_3\}$ & Step 1: $Q_3$ avoids $y_2$; used in Step 3 to exclude $Q_3\ell_3$ as the origin of
$Q_2$ & the origin of $Q_2$ has an end at $y_2$ (Lemma~\ref{lem:fanend}), and
$P_{y_2}\ne P_{y_3}$\\
$\{x_2,y_2\}$ & Step 3: the origin of $Q_2$ avoids $x_2$ & $P_{y_2}$ avoids $x_2$\\
$\{\ell_1,y_2\}$ & Steps 3--4: $Q_2$ and the $\ell_2$-path avoid $\ell_1$ & $Q_1$ does not end at
$y_2$, so $Q_1\ne P_{y_2}$\\
\bottomrule
\end{tabular}
\end{center}
The one new argument is in Step 3. Since $y_2\notin V(P_2)$, $y_2$ is not a target, and since
$y_2\notin N(T_2)$, it is not a chain vertex of any Fan step of the first walk. So
Lemma~\ref{lem:fanend}, applied at each step with $v=y_2$, gives $V(Q_2)\subseteq V(Q)\cup T_2$
for the path $Q$ ending at $y_2$ before the walk. Step 1 extended only $Q_3=P_{y_3}\ne P_{y_2}$,
so $Q=P_{y_2}$ is a path of the base decomposition. It avoids $x_2$ and $\ell_3$, and it avoids
$\ell_1$ because $Q_1$, the only path through the leaf $\ell_1$, does not end at $y_2$. As
$x_2=q_s\notin T_2$ and leaves are not targets, $Q_2$ avoids $x_2$, $\ell_1$ and $\ell_3$, which is
what Steps 3--5 of \cite[Thm~8.2]{TM} use.
\end{proof}

This removes the separation $\{x_2,y_2\}$, and so the obstruction $x_2y_2\in D$, but not
$y_2\notin V(P_2)$. Indeed, $y_2\in V(P_2)$ makes $y_2$ a target, and its new end path is then
created by the Fan step at $y_2$. In particular, if $P_2$ ends with the edge $y_2x_2$ and that
step is an append, the new path ending at $y_2$ is $P_{x_2}$ extended by $y_2$, which contains
$x_2$. So a walk ending with $y_2x_2$ cannot be repaired by tracking alone. Since $q_0=y_3$
is always a target, the condition $y_2\notin N(T_2)$ forces $y_2y_3\in D$, and
\cite[Def.~8.1(ii)]{TM} forces $y_2x_3\in D$. So Proposition~\ref{prop:ep} reaches only shapes in
which $D$ contains the path $x_3\,y_2\,y_3$. In the dense regime
the walks of length one of Section~\ref{sec:main} are therefore the right specialisation: there
every end path is explicit, and no non-adjacency at all is needed beyond the two chain edges.

\section{The base conjecture}\label{sec:conj}

\begin{conjecture}[base]\label{conj:base}
Let $J$ be $3$-connected, $z$ its only even vertex, $d(z)=6$, and
$D=K_6-(J-z)[N(z)]\in\{\emptyset,K_2,P_3,K_3,C_4\}$. Then for every $r\in N(z)$ some labelling with
$r\in\{y_1,y_2,y_3\}$ and $x_1y_2,x_2y_3\in E(J)$ has an append base.
\end{conjecture}

\emph{Conjecture~\ref{conj:base} and T(6)-i.} Suppose Conjecture~\ref{conj:base} holds, and
let $r$ be a root with a neighbour in the core. If the test failed at $r$, the root-shape
theorem~\cite{RS} would put $D$ among the five shapes, and then Corollary~\ref{cor:test1}(a)
shows that $r$ passes. So every root with a core neighbour passes. By \cite[Thm~144]{MR} it
suffices to test roots with at least three neighbours in the core, so T(6)-i follows. A root isolated
in $C$ has $D\supseteq K_{1,5}$ and is not covered by Conjecture~\ref{conj:base} as stated. One
could instead extend the conjecture to $D=K_{1,5}\cup D'$ with $D'\in\{\emptyset,K_2,P_3\}$ and
labellings with $y_1=r$. Labellings with chain edges exist there by the argument of the
$y_1$-case of Lemma~\ref{lem:labelling}, since $D'$ has at most two edges.
Corollary~\ref{cor:test1}(a) is per root, so only the roots not already closed by the known
reductions need a base.

\emph{What kind of statement it is.} By Proposition~\ref{prop:converse} and
Remark~\ref{rem:ptm}, the (E7) form of Conjecture~\ref{conj:base} asks, for each root, for a
labelled pairing $F$ with STL-P and, in addition, two prescribed terminal edges. So it implies
STL-P for those pairings, and it is implied by nothing weaker that we know. It is a statement in
Lov\'asz's setting (one even vertex), like $\STL$ itself, in the family of the separation and
avoidance statements of \cite{MS,L1S}. In the base form each condition (E1)--(E5), (E7) says that
the path ending at a prescribed vertex avoids one or two prescribed vertices. These are exactly
the conditions that undoing the appends produces in the proof of
Proposition~\ref{prop:converse}. The difficulty is that there are six of them, at six vertices of
a nearly complete $K_6$, simultaneously. The known single-vertex tools
(\texttt{lovasz\_avoid\_set}, \texttt{separation\_sparse}) need the special vertices to be packed
or sparse, which is the opposite of the dense regime.

\subsection*{The base conjecture is not a conjunction of $L_1^*$ instances}
Each of (E1)--(E5), (E7) has the form ``the path ending at $c$ avoids $A$'', or can be
strengthened to it. (E1) avoids $x_1$, and $P_{y_2},P_{y_3}$ avoid $z$; for (E4), $P_{y_2}$
avoids $y_3$. The avoidance conjecture $L_1^*$ of \cite{L1S} (open; reduced to TR and ST in the
l1-induction work) treats one such condition in a graph with one even vertex $g$, but only for
$A\cap(N[c]\cup N(g))=\emptyset$. That hypothesis cannot be weakened to $A\cap N(g)\ne\emptyset$:
the example \texttt{DTw} of \cite{L1S} is a counterexample. Here $g=x_3$, and in a dense shape
$x_3$ is adjacent in $B$ to at least three of the other five vertices of $U$.
\texttt{admissible.py} checks all labellings with $x_1y_2,x_2y_3\in E(B)$. In every labelling of
every shape, some conditions avoid a neighbour of $g$. The fewest such conditions are:
\begin{itemize}
\item $C_4$: (E1) and (E7);
\item $K_3$: (E1) and (E2);
\item $P_3$: (E2), (E4) and (E7);
\item $K_2$: (E1), (E2), (E4) and (E7);
\item $\emptyset$: all of (E1)--(E7) except the avoidance of $z$.
\end{itemize}
These conditions are not implied by the others. That is expected from
Proposition~\ref{prop:converse}: each of them is produced when the appends are undone, so dropping
one changes the statement. In residual $C_4$ graphs (\texttt{partial.py}),
$1{,}920$ Lov\'asz decompositions of $M$ were found satisfying only the $L_1^*$-admissible
conditions. Of these, (E1) holds in $186$, (E7) in $1{,}042$, and both in only $120$. So
Conjecture~\ref{conj:base} does not reduce to $L_1^*$, or to TR and ST, constraint by
constraint. Even for $C_4$ it contains an avoidance in the range where $L_1^*$ fails in
general: the $z$-path must avoid $x_1\in N(x_3)$. It is a \emph{joint} statement in the
tight-Fan regime of \cite{TF}. Replacing (E7) by (E6), a predecessor condition, moves one
constraint into the forced-chain form.

\section{Computations}\label{sec:comp}
\emph{This section is computational; no proof above depends on it.} Scripts are in
\texttt{scripts/}; the experiment log is E28 in \texttt{research/experiments.md}.

\subsection*{Method}
\texttt{lsearch.c} is the depth-first search of the star-terminal oracle~\cite{STL,TM}: paths are
grown from a vertex with an end, with failures memoised by the remaining edges and end profile.
It searches for path decompositions with prescribed end counts under constraints that depend only
on the path being closed: end-path avoidance, forbidden end pairs, the end on a given side of a
given edge, and separation. So the memoisation stays exact. A node budget can stop a search
(``undecided'').

\texttt{dense.py} loops over all roots $r$ and all labellings with $r\in\{y_i\}$ and
$x_1y_2,x_2y_3\in E(B)$, in random order. It searches for a base and then executes the whole
construction on the base found.
\begin{itemize}
\item Mode \texttt{lean} is the exact hypothesis \texttt{AppendBase}. It tries (E7) first and
otherwise (E6), and with (E6) it also checks that the static condition holds in $\cD_3$.
\item Mode \texttt{endZ} is (E1)--(E5) with a strengthened (E6): every leaf is identified with
$z$, so the end on $p$'s side must also avoid $z$ and $y_3$. Its bases satisfy (E6).
\item Mode \texttt{endA} is (E1)--(E5) with (E6) minus the clause $w=y_3$. That clause is
redundant (A26, \S1.2: the transported predecessor read from $y_3$ never arises, because $y_3$ has
no end in $\cD_1$; Remark~\ref{rem:e6}). The Lean theorem uses the full (E6), so for
\texttt{endA} bases of \texttt{dense.py} the guarantee is on paper only; they are also executed. (In \texttt{c4.py}, mode \texttt{endA}
uses the $C_4$ labelling, where $x_2y_2,y_2y_3\notin E(B)$, and is exactly (E1)--(E6).)
\item Mode \texttt{end} is (E1)--(E4) followed by Fan walks; it is the hypothesis of no theorem
here.
\item Mode \texttt{six-separation} (\texttt{c4.py sep}) is the base of \cite[Thm~8.2]{TM} with the
chains $P_2=y_3x_2$, $P_1=y_2x_1$. For $C_4$ these violate \cite[Def.~8.1(ii)]{TM}, so the
protected step is executed, not guaranteed.
\end{itemize}
The construction is run by \texttt{walksim.py} with forced Fan chains (on decompositions with at
most one end per vertex, Fan's chain is forced: \texttt{fanPred\_unique}). The final
decomposition of $J$ is then checked edge by edge: every path simple, every edge of $J$ covered
exactly once, end counts, and the star condition.

\subsection*{Graphs}
\begin{itemize}
\item \texttt{gen\_core.py} \cite{TM}: random $3$-connected $J$ with $B[U]=K_6-D$, exterior
degrees $3$ or $5$, one or three exterior neighbours per vertex of $U$; a sparse exterior.
\item \texttt{gen\_resid.py}: random $3$-connected $J$ with $B[U]=K_6-D$ and some of the
residual conditions of a smallest failure imposed:
  \begin{itemize}
  \item $\delta(B)\ge6$ and $\delta(B[O])\ge4$;
  \item every vertex of $U$ has at least three odd neighbours, a universal one at least five;
  \item $|N(x)\cap U|\in\{0,3,4,5,6\}$ for $x\in O$;
  \item for $C_4$, the two pairs of opposite cycle vertices are closed twins in $B$ (so every
  cycle vertex qualifies as a failing root), or one pair is (TW$=1$);
  \item for $\emptyset$, an exterior vertex adjacent to all of $U$.
  \end{itemize}
  With $|O|=8$ these graphs have $|B|=14$ and $64$--$82$ edges. Not imposed: twins of degree at
  least eight for $C_4$, $d(r)\ge8$ with three core neighbours \cite[Thm~144]{MR}, and the
  bad-vertex table of the reader notes.
\end{itemize}

\subsection*{Results}
Every root of every graph is tested, not only the roots allowed by the shape theorem. A root
counts as certified if some labelling has a base and the executed construction verifies.

\emph{Exhaustive, order $9$.} Among the $1{,}819$ $3$-connected graphs $J$ of order $9$ in which
$z$ is the only even vertex and $d(z)=6$ (the list of \cite[\S8]{TM}), $28$ have a dense core.
They split as $2$ $\emptyset$, $6$ $K_2$, $8$ $P_3$, $5$ $K_3$ and $7$ $C_4$. With exact search
(no node budget) in mode \texttt{lean}:
\begin{itemize}
\item every one of the $168$ roots is certified at the first labelling tried, through (E7);
\item of all $42{,}432$ pairs (root, labelling with $x_1y_2,x_2y_3\in E(B)$), $42{,}384$ have a
base with (E7) and $24$ more with (E6);
\item $24$, all of shape $C_4$, have none;
\item every root has a base for at least $96.7\%$ of its labellings;
\item no executed construction failed.
\end{itemize}

\emph{Exhaustive, order $9$: the cost of the strengthening.} For each of the $14{,}144$
labellings with chain edges we also decided STL-P for the oriented pairing (a star-terminal
Lov\'asz decomposition of $J$ with transitions $x_i\,z\,y_i$, $x_i$ an end; exact search,
\texttt{stlp.py}; side (b) of Proposition~\ref{prop:converse} is \texttt{equiv.py}).
STL-P holds for all $14{,}144$. An (E7)-base exists for $14{,}128$, in agreement with
Proposition~\ref{prop:converse}: the auditor's search of side (b) finds the same $14{,}128$. An
(E6)-base exists for $8$ of the remaining $16$. The other $8$ labellings, all of shape $C_4$,
have no base of either kind although STL-P holds for them. So the strengthening is strict on
single labellings, but at order $9$ it costs no root.

\emph{Random, residual-style and sparse} (node budget $2\cdot10^6$ to $5\cdot10^6$ per search,
two to four seeds). ``Failures'' counts bases satisfying the mode's conditions whose executed
construction did not verify.
\begin{center}\footnotesize
\begin{tabular}{@{}llrrrr@{}}
\toprule
graphs & mode & $J$ & roots & certified & failures\\
\midrule
\texttt{gen\_core} $C_4$, $|B|=14$ & six-separation & 10 & 40 & 40 & 216 of 1{,}050\\
\texttt{gen\_core} $C_4$, $|B|=14$ & end & 10 & 40 & 40 & 288 of 1{,}280\\
\texttt{gen\_core} $C_4$, $|B|=14$ & endA & 10 & 40 & 40 & 0 of 1{,}280\\
\texttt{gen\_core} $\emptyset,K_3,C_4$, $|B|=14$ & endA & 18 & 108 & 108 & 0\\
residual, 3 per shape, $|O|=8$ & endA and endZ & 15 & 90 & 90 & 0\\
residual $C_4$ (twins), $|O|=8$ & endZ & 20 & 120 & 120 & 0\\
residual $K_3$, $|O|=8$ & endZ & 20 & 120 & 120 & 0\\
residual $P_3$, $|O|=8$ & endZ & 20 & 120 & 114 & 0\\
residual $K_2$, $|O|=8$ & endZ & 4 & 24 & 24 & 0\\
residual $K_2$, $C_4$, $P_3$ (from the 3 per shape) & lean & 8 & 48 & 48 & 0\\
residual $C_4$ (twins), $|O|=8$ & lean & 8 & 48 & 48 & 0\\
residual $K_3$, $|O|=8$ & lean & 8 & 48 & 48 & 0\\
residual $P_3$, $|O|=8$ & lean & 8 & 48 & 48 & 0\\
residual $\emptyset$, $|O|=8$ & lean${}^*$ & 8 & 48 & 42 & 0\\
residual $K_2$, $|O|=8$ & lean${}^*$ & 8 & 48 & 30 & 0\\
residual $C_4$ (twins), $|O|=10$ & lean${}^*$ & 8 & 48 & 48 & 0\\
residual $K_3$, $|O|=10$ & lean${}^*$ & 8 & 48 & 42 & 0\\
residual $P_3$, $|O|=10$ & lean${}^*$ & 8 & 48 & 44 & 0\\
residual $\emptyset$, $|O|=10$ & lean${}^*$ & 8 & 48 & 48 & 0\\
residual $C_4$ (one twin pair), $|O|=10$ & lean${}^*$ & 8 & 48 & 48 & 0\\
\bottomrule
\end{tabular}
\end{center}
(${}^*$: at most $60$ labellings per root.) The six-separation and \texttt{end} bases are not
hypotheses of a theorem here. Their failures are protected steps that hit a blank, which is the
obstruction of \cite[\S9(1)]{TM}. No base satisfying (E1)--(E5) with (E6) or (E7), in any mode,
failed to verify. In mode \texttt{lean} almost every certificate uses (E7), so that all seven
steps are appends. (E6) was needed for $12$ of the $536$ certified roots in mode \texttt{lean}.

\emph{Uncertified roots.} They lie in eight graphs, $46$ roots in all: $6$ in the
\texttt{endZ} $P_3$ row; $4$, $6$ and $6$ in the \texttt{lean} rows $P_3$, $K_3$ and $K_2$ with
$|O|=10$; $6$ in the $\emptyset$ row and $18$ (three graphs) in the $K_2$ row with $|O|=8$. In
each of these graphs every search of every tried labelling hit the node budget. For the five
graphs found last this was rechecked with a count of proved absences, and it is zero. In each
such graph all or most roots are uncertified, which points at the search rather than at the
graph. On the first of the hard $K_2$ graphs, the weaker statement STL-P for the labelled pairing
is also undecided at the same budget for $20$ of $20$ labellings. So these roots are undecided,
not counterexamples.

For comparison, the six-separation base of \cite[Thm~8.2]{TM}, with the same chains, was tried on
five residual $C_4$ graphs ($|O|=8$; all $160$ labellings with a root on the cycle; two seeds;
budget $5\cdot10^6$). Every search hit the budget. So no base was found, but its absence is not
proved. On the sparse \texttt{gen\_core} graphs the six-separation base exists for $304$ of the
$320$ configurations, and the protected step then succeeds for $834$ of the $1{,}050$ bases found
(forced chain).

\section{Relation to the known reductions of T(6)-i}\label{sec:novelty}

The known closures and reductions of T(6)-i (STATUS~\S3.5(a); reader notes; audits A6, A9) are
the following.
\begin{itemize}
\item The root-shape theorem: $D\in\{\emptyset,K_2,P_3,K_3,C_4\}$, with the root restrictions.
\item Six-even-incidence Thm~1: odd vertices with one or two even neighbours.
\item The four-exterior paper: single-witness closures by degree.
\item The five-odd paper.
\item The six-clique theorem: $D=\emptyset$ without full exterior vertex.
\item The twin results for $C_4$: adjacent-twin insertion at degree six, partial at degree eight.
\item The degree-ten note.
\item MR Thm~144 and Lemma~133.
\end{itemize}
They close configurations by conditions on the exterior: contact sizes, degrees, twins. Beside
them there are reformulations, several with one even vertex:
\begin{itemize}
\item $\STL$ and \texttt{test1\_of\_stl} \cite{STL} (synthesis~\S19): Test1 from star-terminal
Lov\'asz decompositions of $J$;
\item STL-P, prescribing the pairing at $z$, and its equivalence with the terminal-matching
property PTM$_3$ of $(J-z)+F$ (synthesis~\S20; \cite{TM});
\item the chained walk theorem \cite[Thm~8.2]{TM}, which reduces PTM$_3$ to a base with one even
vertex and six separations;
\item $B_3\Leftrightarrow$ T(6)-i (endpoint-structure Cor.~30, \texttt{rooted\_six\_of\_bstmt3});
\item the hub transfers and certificates of the d6-frontier paper (audit A9).
\end{itemize}
The dense shapes remain open in general:
\begin{itemize}
\item $K_2$, $P_3$, $K_3$ have no shape-specific result;
\item $C_4$ is open for twins of degree at least eight;
\item $\emptyset$ is open with full exterior vertices.
\end{itemize}
The present paper does \emph{not} close any of these. It reformulates all five shapes,
uniformly and without exterior hypotheses, as Conjecture~\ref{conj:base}.

\emph{Position in the chain of reformulations.} We have
\[
\text{Conj.~\ref{conj:base}}\ \Longrightarrow\ \text{STL-P (labelled pairings)}\ \Longrightarrow\
\STL\ \Longrightarrow\ \mathrm{Test1}.
\]
By Proposition~\ref{prop:converse}, the (E7) form of Conjecture~\ref{conj:base} is equivalent to
STL-P with two further prescribed terminal edges. The first implication is strict on single
labellings: at order $9$ some labellings satisfy STL-P but have no base
(Section~\ref{sec:comp}). Nothing here shows that Conjecture~\ref{conj:base} is easier than
STL-P. What the reformulation gains is a base statement on the fixed graph $M$ that names, at each
step of the construction, the path that must avoid a vertex, with no chain hypothesis.

\emph{Comparison with the composition method.} The composition method of \cite{TM} produced a
new proof only of a known case: far-apart neighbours, covered by six-even-incidence Thm~1(2)
\cite[Rem.~8.5]{TM}. Its base hypothesis is unsatisfiable for $D\ne C_4$, and with
\cite[Def.~8.1(ii)]{TM} the theorem as stated applies to no dense shape; both observations are
from \cite[\S9(1)]{TM} and connections~\S Z.4. End-path tracking, proposed there, is carried out
here for walks of length one. In the computed dense instances the base exists for almost every
root, for all five shapes (Section~\ref{sec:comp}).

\section{The formalisation}\label{sec:lean}
\sloppy
The file \texttt{T6Dense.lean} in \texttt{lean/GallaiLean} contains no \texttt{sorry}. The
file \texttt{T6DenseAxioms.lean} in this directory prints only \texttt{propext},
\texttt{Classical.choice} and \texttt{Quot.sound} for the main theorems, and the full library
builds.
\begin{itemize}
\item \texttt{delList}, \texttt{delList\_cons}, \texttt{delList\_single}: a graph minus a list of
edges.
\item \texttt{append\_full}: the append move with the new path's vertex sequence;
\texttt{SameP}, \texttt{ExtP} and their lemmas: bookkeeping of paths through appends;
\texttt{fanPred\_append}, \texttt{fanPred\_self}: Lemma~\ref{lem:transport}.
\item \texttt{splitG}, \texttt{spl}, \texttt{splBase}: the split graph and the graphs of the
construction; \texttt{collapse} with \texttt{collapse\_hE}, \texttt{collapse\_hEinj},
\texttt{collapse\_hcov}: the identification.
\item \texttt{split\_abc} (steps 1--3), \texttt{LeafOK}, \texttt{leafOK\_three},
\texttt{LeafOK.of\_cover}, \texttt{split\_finish} (steps 5--8), \texttt{split\_stage4}.
\item \texttt{star\_of\_appendBase} (Theorem~\ref{thm:append}), \texttt{star\_of\_protectedBase}
(Theorem~\ref{thm:protected}, from \texttt{walk\_step\_protected} of \texttt{TightFan.lean}),
\texttt{AppendBase}, \texttt{stl\_of\_appendBase}, \texttt{test1\_of\_appendBases}
(Corollary~\ref{cor:test1}, from \texttt{test1\_of\_stl}).
\item \texttt{baseM}, \texttt{baseM\_parity}, \texttt{LovaszBase},
\texttt{appendBase\_of\_lovaszBase}, \texttt{stl\_of\_lovaszBase},
\texttt{test1\_of\_lovaszBases}: the same with the hypothesis stated on $M$ itself, as in
Definition~\ref{def:base}.
\end{itemize}
The Lean conclusions of \texttt{star\_of\_appendBase} and \texttt{star\_of\_protectedBase} are
the $\STL$ data (star-terminal at $z$, and a long side at $y_i$ or a path of length two), not the
transitions $x_i\,z\,y_i$. Lemmas~\ref{lem:labelling} and~\ref{lem:fanend},
Propositions~\ref{prop:converse} and~\ref{prop:ep}, Corollary~\ref{cor:test1}(a) as a separate
statement, and the step from Conjecture~\ref{conj:base} to T(6)-i through \cite{RS} and
\cite[Thm~144]{MR} are not formalised.

\begin{thebibliography}{9}
\bibitem{Fan05} G.~Fan, Path decompositions and Gallai's conjecture, J.~Combin.\ Theory Ser.~B 93
(2005) 117--125.
\bibitem{INC} Six-even-incidence (Nov.~2026 study of the collection), Theorem~1; summarised in
\texttt{research/notes/reader-six-even-program.md}, audited in A6.
\bibitem{L1S} J.~Williams, The avoidance conjecture $L_1^*$, \texttt{manuscripts/l1star}, 2026.
\bibitem{MR} J.~Williams, Marked reduction, Nov.~2026 (Thms~114, 123, 133, 144).
\bibitem{MS} J.~Williams, Separating many pairs in Lov\'asz decompositions,
\texttt{manuscripts/multi-sep}, 2026.
\bibitem{RS} Six-even-root-shapes (Nov.~2026 study of the collection); summarised in the reader
notes.
\bibitem{SP} J.~Williams, The spine: Gallai's conjecture and two deletion tests on $3$-connected
graphs, \texttt{manuscripts/spine}, 2026.
\bibitem{STL} J.~Williams, Star-terminal Lov\'asz decompositions and the deletion tests,
\texttt{manuscripts/stl}, 2026.
\bibitem{TF} J.~Williams, Fan's lemma at equality, \texttt{manuscripts/tight-fan}, 2026.
\bibitem{TM} J.~Williams, Two terminal edges at once: the terminal-matching property of Lov\'asz
decompositions, \texttt{manuscripts/terminal-matching}, 2026.
\end{thebibliography}
\end{document}
